\section{Coefficient reconstruction for boundary rays}
\label{app:three-boundary}

This appendix proves the rank statement used in
\Cref{three:boundary} by coefficient elimination. No generic-position
assumption or computational rank test is needed.

The family has the following exact edge restrictions, where
$D=d_1+d_2-h$:
\[
\begin{aligned}
 q(x,0,0)&=(h-d_1x)^2,&q(0,y,0)&=(h-d_2y)^2,\\
 q(1,0,z)&=(d_3z-(d_1-h))^2,&
 q(0,1,z)&=(d_3z-(d_2-h))^2,\\
 q(1,1,z)&=(d_3z-(D+k))^2.
\end{aligned}
\]
These identities verify every edge value and tangency used below.
At $h=0$, the origin additionally has zero $x$- and $y$-derivatives.
At $h=d_1$, the vertex $(1,0,0)$ has zero $x$- and $z$-derivatives.

Write a candidate quadratic as
\begin{equation}\label{eq:three:candidate}
 p=A+a x+b y+c z+u x^2+v y^2+w z^2+rxy+sxz+tyz.
\end{equation}
Let
\[
 \gamma=\frac{d_1-h}{d_3},\qquad
 \delta=\frac{d_2-h}{d_3},\qquad
 \eta=\frac{d_1+d_2-h+k}{d_3}.
\]
In every regime under consideration, the three $z$-edges at $(1,0)$,
$(0,1)$, and $(1,1)$ carry a double zero at $\gamma$, $\delta$, and $\eta$,
respectively. The term ``double zero'' means that both the value and the
edge derivative vanish; it includes an endpoint with the stated
tangency. The corresponding restrictions are therefore
\begin{align}
 p(1,0,z)&=w(z-\gamma)^2,\label{eq:three:edge-c}\\
 p(0,1,z)&=w(z-\delta)^2,\label{eq:three:edge-d}\\
 p(1,1,z)&=w(z-\eta)^2.\label{eq:three:edge-e}
\end{align}
These are identities even if $w=0$. Their linear coefficients give
\begin{equation}\label{eq:three:z-coefficients}
 c=2w(\eta-\gamma-\delta)=\frac{2w(h+k)}{d_3},
 \quad s=-\frac{2w(d_1+k)}{d_3},
 \quad t=-\frac{2w(d_2+k)}{d_3}.
\end{equation}
Their constant coefficients give
\begin{equation}\label{eq:three:edge-constants}
 A+a+u=w\gamma^2,\qquad A+b+v=w\delta^2,
 \qquad A+a+b+u+v+r=w\eta^2.
\end{equation}

\subsection{The upper contact reaches the vertex}

Suppose $0<h<\min(d_1,d_2)$ and
$d_3=d_1+d_2-h+k$, so $\eta=1$. The additional $x$- and $y$-edge
contacts at $h/d_1$ and $h/d_2$ are interior double zeros. Thus
\[
 A=u\frac{h^2}{d_1^2}=v\frac{h^2}{d_2^2},\qquad
 a=-\frac{2uh}{d_1},\qquad b=-\frac{2vh}{d_2}.
\]
The first equation in \eqref{eq:three:edge-constants} becomes
\[
 \frac{A(d_1-h)^2}{h^2}
       =\frac{w(d_1-h)^2}{d_3^2}.
\]
Since $d_1>h$, cancellation is valid. Hence
\[
 A=\frac{wh^2}{d_3^2},\quad
 u=\frac{wd_1^2}{d_3^2},\quad
 v=\frac{wd_2^2}{d_3^2},\quad
 a=-\frac{2whd_1}{d_3^2},\quad
 b=-\frac{2whd_2}{d_3^2}.
\]
Together with \eqref{eq:three:z-coefficients}, the final constant equation
determines $r$. All ten coefficients are fixed by $w$.

\subsection{The two lower contacts reach the origin}

Suppose $h=0$ and $d_3\ge d_1+d_2+k$. At the origin, the value and
the $x$- and $y$-derivatives are zero, so $A=a=b=0$.
The first two equations in \eqref{eq:three:edge-constants} give
\[
 u=\frac{wd_1^2}{d_3^2},\qquad
 v=\frac{wd_2^2}{d_3^2}.
\]
Equations \eqref{eq:three:z-coefficients} and
\eqref{eq:three:edge-constants} determine the remaining coefficients.
There is no division by $d_1$ or $d_2$, so this proof includes either or
both of them equal to zero. In that event, a zero on a $z$-edge reaches
its lower vertex and its $z$-derivative remains zero. If
$d_3=d_1+d_2+k$, the upper contact reaches $(1,1,1)$ with its
$z$-derivative zero. Thus the same equations apply at both ends of this
parameter regime.

\subsection{Two contacts merge at a lower vertex}

Suppose $h=d_1<d_2$ and $d_3\ge d_2+k$. The contact at $(1,0,0)$
has zero value and zero $x$- and $z$-derivatives. Consequently,
\[
 A=u,\qquad a=-2u,\qquad c+s=0.
\]
The $y$-edge contact is at $\beta=d_1/d_2<1$ and has zero derivative.
For $d_1=0$, this is the origin with a $y$-tangency. In both cases,
\[
 A=v\beta^2,\qquad b=-2v\beta.
\]
The second equation of \eqref{eq:three:edge-constants} becomes
\[
 v(1-\beta)^2=\frac{w(d_2-d_1)^2}{d_3^2}.
\]
Since $d_2>d_1$, it follows that
\[
 v=\frac{wd_2^2}{d_3^2},\quad
 A=u=\frac{wd_1^2}{d_3^2},\quad
 a=-\frac{2wd_1^2}{d_3^2},\quad
 b=-\frac{2wd_1d_2}{d_3^2}.
\]
The formulas for $c,s,t$ and $r$ follow as before. This argument also
covers $d_1=0$, because it divides only by $d_2-d_1$ and $d_2>0$.
The upper contact can be either an interior point of the $z$-edge or its
upper vertex, depending on whether $d_3>d_2+k$ or equality holds.

\subsection{The common one-dimensional kernel}

In all three regimes, the resulting vector of coefficients is
\begin{align*}
 (A,a,b,c,u,v,w,r,s,t)=\frac{w}{d_3^2}\bigl(&h^2,-2hd_1,-2hd_2,
       2d_3(h+k),d_1^2,d_2^2,d_3^2,\\
       &2d_1d_2+k(2(d_1+d_2-h)+k),
       -2d_3(d_1+k),-2d_3(d_2+k)\bigr).
\end{align*}
This is exactly $(w/d_3^2)q_{h,d,k}$. Since $d_3>0$ in every stated
regime, the family member is nonzero and supplies a kernel vector.
Every kernel vector has the displayed form, including the zero vector
when $w=0$. The contact system therefore has rank nine in the
ten-dimensional quadratic space throughout each regime. This proves
the claim used in \Cref{three:boundary}.
